\chapter{冷热分层存储对照实验}
\label{ch:hotcold}

\section{实验目的与方案设置}

第四章验证的是基础多层级结构中 facility、k-kick、局部路由和重构模块能否正常工作。第三章已经说明冷热分层键存储的实现方式：插入的键先进入内存 tail cubby，tail cubby 容量满，切换时再批量下沉到 SQLite3 数据库存储。本章在该实现基础上设置三组对照实验，观察内存 tail cubby 与 SQLite3 冷数据层结合后带来的系统性能效果。第五章与第四章延迟数值不可直接对比，是因为测试的数据方式不同，属于合理范围。

实验设置了三个对照方案：(1) 原始的多层级 cubby 结构，数据均保存在内存多层级 cubby 中；(2) 纯数据库结构，所有插入、查询和删除都直接访问 SQLite3；(3) 冷热分层结构，近期写入的数据先进入内存 tail cubby，tail cubby 写满后再批量写入 SQLite3。在相同的数据规模和相同的测试数据情况下进行实验，避免数据层面的差异带来的误差。

\section{负载与统计口径}

实验选取 \(5\times10^3\)、\(10^4\)、\(5\times10^4\)、\(10^5\)、\(5\times10^5\) 五档规模。每一档先插入全部键，再进行同规模次数的查询，随后删除近期的 20\% 键。测试数据采用 C++ 标准库的 std::mt19937\_64 随机数生成器产生互不重复的 64 位整数键，以保证键值分布的均匀性和实验结果的可复现性。完全插入系统后，查询阶段按照键的插入顺序建立访问排名，其中最新插入的键具有更高访问概率，并依据 Zipf 分布进行采样，以模拟具有时间局部性的访问负载\cite{cooper2010ycsb}。删除阶段选择最近插入的 20\% 键，用于模拟查询与删除主要针对近期写入的键。

统计指标包括插入平均延迟、查询平均延迟、删除平均延迟、平均每键空间占用和正确率。平均延迟只统计对应操作本身，不把收尾阶段的数据库刷盘和检查过程计入单次操作。平均每键空间占用按删除后的剩余键数计算，内存结构统计逻辑元数据位数，数据库结构统计 SQLite3 主数据库文件大小。

\section{多层级内存结构对照}

表~\ref{tab:hotcold-memory-latency} 对比了多层级内存结构与冷热分层结构的实验结果。多层级内存结构的查询延迟在五档规模下均低于 1\,$\mu$s，说明第四章中的局部路由在新增负载下仍然保持较短查询路径。冷热分层结构的查询延迟为 0.399--0.550\,$\mu$s，除小规模 \(5\times10^3\) 外，与多层级内存结构处于同一数量级。

\begin{table}[htbp]
  \centering
  \caption{多层级内存结构与冷热分层结构延迟对照}
  \label{tab:hotcold-memory-latency}
  \small
  \begin{tabular}{llrrr}
    \toprule
    规模 & 方案 & 插入 ($\mu$s) & 查询 ($\mu$s) & 删除 ($\mu$s) \\
    \midrule
    \(5\times10^3\)   & 多层级内存结构 & 42.189 & 0.335 & 140.200 \\
    \(5\times10^3\)   & 冷热分层结构   & 10.231 & 0.476 & 8.605 \\
    \(10^4\)          & 多层级内存结构 & 41.804 & 0.345 & 123.759 \\
    \(10^4\)          & 冷热分层结构   & 9.686  & 0.399 & 11.341 \\
    \(5\times10^4\)   & 多层级内存结构 & 57.381 & 0.510 & 181.186 \\
    \(5\times10^4\)   & 冷热分层结构   & 17.198 & 0.464 & 11.148 \\
    \(10^5\)          & 多层级内存结构 & 61.702 & 0.516 & 167.669 \\
    \(10^5\)          & 冷热分层结构   & 17.432 & 0.486 & 12.638 \\
    \(5\times10^5\)   & 多层级内存结构 & 83.232 & 0.708 & 129.434 \\
    \(5\times10^5\)   & 冷热分层结构   & 19.271 & 0.550 & 23.660 \\
    \bottomrule
  \end{tabular}
\end{table}

插入和删除结果更能体现两类结构的差异。多层级内存结构需要维护 tail cubby 回填以及层级合并、拆分，删除延迟在五档规模下为 123.759--181.186\,$\mu$s，\(5\times10^3\) 时为 140.200\,$\mu$s。冷热分层结构删除近期键时多发生在内存 tail cubby 或尚未完全下沉的数据区域，删除延迟为 8.605--23.660\,$\mu$s。在 Zipf 写后读负载下，冷热分层插入和删除延迟更低，删除延迟比多层级内存结构低一个数量级。

\section{纯数据库结构对照}

表~\ref{tab:hotcold-sqlite-latency} 对比了纯数据库结构与冷热分层结构的实验结果。纯数据库结构的查询都需要经过 SQLite3\cite{gaffney2022sqlite}，五档规模下查询延迟为 1.190--1.433\,$\mu$s；冷热分层结构在相同查询序列下为 0.399--0.550\,$\mu$s。该差异来自负载中的近期访问特征：一部分新近写入的键仍在内存 tail cubby 中，查询不必进入数据库。

\begin{table}[htbp]
  \centering
  \caption{纯数据库结构与冷热分层结构延迟对照}
  \label{tab:hotcold-sqlite-latency}
  \small
  \begin{tabular}{llrrr}
    \toprule
    规模 & 方案 & 插入 ($\mu$s) & 查询 ($\mu$s) & 删除 ($\mu$s) \\
    \midrule
    \(5\times10^3\)   & 纯数据库结构 & 12.064 & 1.190 & 13.575 \\
    \(5\times10^3\)   & 冷热分层结构 & 10.231 & 0.476 & 8.605 \\
    \(10^4\)          & 纯数据库结构 & 12.639 & 1.257 & 13.735 \\
    \(10^4\)          & 冷热分层结构 & 9.686  & 0.399 & 11.341 \\
    \(5\times10^4\)   & 纯数据库结构 & 13.659 & 1.279 & 12.301 \\
    \(5\times10^4\)   & 冷热分层结构 & 17.198 & 0.464 & 11.148 \\
    \(10^5\)          & 纯数据库结构 & 13.684 & 1.284 & 13.464 \\
    \(10^5\)          & 冷热分层结构 & 17.432 & 0.486 & 12.638 \\
    \(5\times10^5\)   & 纯数据库结构 & 20.567 & 1.433 & 25.888 \\
    \(5\times10^5\)   & 冷热分层结构 & 19.271 & 0.550 & 23.660 \\
    \bottomrule
  \end{tabular}
\end{table}

插入延迟并不是冷热分层结构在所有规模下都更低，\(5\times10^4\) 与 \(10^5\) 两档中，冷热分层结构插入分别为 17.198\,$\mu$s 和 17.432\,$\mu$s，高于纯数据库结构的 13.659\,$\mu$s 和 13.684\,$\mu$s，但是这处于能接受的范围内。这与内存 tail cubby 中的 k-kick 插入以及后续批量写入有关。到 \(5\times10^5\) 时，冷热分层结构插入为 19.271\,$\mu$s，略低于纯数据库结构的 20.567\,$\mu$s。删除方面，冷热分层结构五档均低于纯数据库结构，但差距不如查询明显。因此，收益主要体现在近期查询与删除，写入并不稳定优于纯 SQLite3。

\section{空间占用与正确性对照}

表~\ref{tab:hotcold-space-correctness} 汇总三种方案的平均每键空间占用和正确率的实验结果。在本文测试用例与访问序列下，三类结构的查询结果均与预期一致，说明原型实现保持了基本插入、查询和删除语义。空间方面，多层级内存结构为 64.427--68.305 位/键，明显低于使用 SQLite3 的方案。这是因为多层级内存结构只统计逻辑元数据位数，而 SQLite3 主数据库文件包含页面、表结构和索引等额外的固定开销。

\begin{table}[htbp]
  \centering
  \caption{三种方案空间占用与正确率对照}
  \label{tab:hotcold-space-correctness}
  \small
  \begin{tabular}{llrr}
    \toprule
    规模 & 方案 & 每键位数 & 正确率（\%） \\
    \midrule
    \(5\times10^3\)   & 多层级内存结构 & 65.333  & 100.000 \\
    \(5\times10^3\)   & 纯数据库结构   & 204.800 & 100.000 \\
    \(5\times10^3\)   & 冷热分层结构   & 188.416 & 100.000 \\
    \(10^4\)          & 多层级内存结构 & 64.427  & 100.000 \\
    \(10^4\)          & 纯数据库结构   & 200.704 & 100.000 \\
    \(10^4\)          & 冷热分层结构   & 184.320 & 100.000 \\
    \(5\times10^4\)   & 多层级内存结构 & 67.141  & 100.000 \\
    \(5\times10^4\)   & 纯数据库结构   & 160.563 & 100.000 \\
    \(5\times10^4\)   & 冷热分层结构   & 154.829 & 100.000 \\
    \(10^5\)          & 多层级内存结构 & 67.904  & 100.000 \\
    \(10^5\)          & 纯数据库结构   & 160.973 & 100.000 \\
    \(10^5\)          & 冷热分层结构   & 153.600 & 100.000 \\
    \(5\times10^5\)   & 多层级内存结构 & 68.305  & 100.000 \\
    \(5\times10^5\)   & 纯数据库结构   & 155.075 & 100.000 \\
    \(5\times10^5\)   & 冷热分层结构   & 151.880 & 100.000 \\
    \bottomrule
  \end{tabular}
\end{table}

在使用 SQLite3 的两类方案之间，冷热分层结构的每键位数略低于纯数据库结构，例如 \(10^5\) 时为 153.600 位/键，而纯数据库结构为 160.973 位/键。这一差异与部分近期数据仍保留在内存 tail cubby、数据库文件中实际写入量和页面分配情况有关。该结果不能简单解释为冷热分层结构空间上总是更优，因为内存 tail cubby 的逻辑元数据和数据库页面开销属于不同来源；但在本实验口径下，冷热分层结构没有因为增加内存热层而超过纯数据库结构的总空间统计，相反因为原本的紧凑哈希设计，存储在 tail cubby 中的数据占用更少的空间。

\section{本章小结}

本章用相同输入和统计口径比较了多层级内存结构、纯数据库结构和冷热分层结构。在具有时间局部性的查询与删除负载下，冷热分层带来的性能收益来自内存热层：近期查询较少直接访问 SQLite3，近期删除也较少触发多层级回填。写入并不总是占优，\(5\times10^4\) 和 \(10^5\) 两档下的插入延迟高于纯数据库结构。若查询目标大多落在冷层，或者 SQLite3 已经有较高页面缓存命中率，热层判断和状态维护反而可能增加成本。因此，本章结论只对应访问局部性较强的负载，不能解释为冷热分层无条件更快。
